\section{Experimental Design}
\label{sec:experiment}

\subsection{Research questions}
\begin{itemize}
    \item \textbf{RQ1 (the labelling function as an analysis).} How accurately
          does $\phi$ identify memory and string operations, measured against
          independent blind annotation on the corpora used in the
          experiments?
    \item \textbf{RQ2 (which node view transfers).} With the graph, backbone
          and training budget fixed, how does cross-project detection change
          when the node view changes? In particular: does the grammar kind
          $\kappa$ beat text; does it beat text \emph{plus} node type; does it
          beat anonymised text; and does the engineered taxonomy $\phi$ add
          anything to $\kappa$?
    \item \textbf{RQ3 (generality).} Does the effect hold across two corpora
          and across four architectures, including a Transformer?
    \item \textbf{RQ4 (ablations).} Which parts of the graph and of the
          taxonomy matter?
    \item \textbf{RQ5 (reference points).} How do the results compare with
          re-implementations of existing input representations, with a
          pretrained code model, and with a trivial function-size rule?
\end{itemize}

\subsection{Corpora, sampling and class balance}
\label{sec:corpora}

We use BigVul~\cite{fan2020bigvul} and DiverseVul~\cite{chen2023diversevul},
loaded from pinned public mirrors, plus Devign~\cite{zhou2019devign} for the
analysis validation only. From each of BigVul and DiverseVul we draw
\numcfgN{} functions \emph{uniformly at random} without replacement (a fixed
seed, so the sample is reproducible), after which near-duplicates are removed
(Section~\ref{sec:dedup}). This replaces an earlier design that took the first
rows in corpus order; it also raises the number of vulnerable functions from
about 430 to \numbigvulVuln{} (BigVul) and \numdiversevulVuln{} (DiverseVul).
Table~\ref{tab:corpora} gives the resulting properties.

% Generated by experiments/emit_tables.py -- DO NOT EDIT BY HAND.
% Every value is read from experiments/results/*.json.

\begin{table*}[t]
\centering
\caption{Corpus properties measured on our uniform random samples
(10{,}000 functions drawn per corpus, then de-duplicated). \emph{Prevalence} is
the fraction of vulnerable functions in the whole sample; \emph{Test prev.} is
the pooled fraction over the held-out folds. They differ because the largest
repositories, which have a lower share of vulnerable functions, are held
train-only and never appear in a test fold (Section~\ref{sec:protocol}).
\emph{Truncated} is the percentage of functions with more than 200 named AST
nodes, overall / vulnerable / benign.}
\label{tab:corpora}
\resizebox{\linewidth}{!}{%
\begin{tabular}{lrrrrrrrr}
\toprule
Corpus & Functions & Vulnerable & Prevalence & Test prev. & Projects & Top-2 share & Nodes (mean) & Truncated (\%) \\
\midrule
BigVul & 9,435 & 555 & 0.059 & 0.071 & 250 & 64.3\% & 96 & 18.1 / 42.3 / 16.6 \\
DiverseVul & 9,791 & 567 & 0.058 & 0.059 & 633 & 25.0\% & 110 & 25.6 / 53.4 / 23.9 \\
Devign & 9,666 & 4,455 & 0.461 & --- & 2 & 100.0\% & 137 & 41.2 / 44.3 / 38.6 \\
\bottomrule
\end{tabular}}
\end{table*}


Table~\ref{tab:node_dist} gives the distribution of operation types over
the nodes of the BigVul sample: about three quarters of nodes carry no
operation label ($\bot$) but keep their grammar kind, and memory operations are
rare.

% Generated by experiments/emit_tables.py -- DO NOT EDIT BY HAND.
% Every value is read from experiments/results/*.json.

\begin{table}[t]
\centering
\caption{Node-type distribution under $\phi_8$ (BigVul sample). $\bot$
(\texttt{UNKNOWN}) nodes carry no operation label but keep their grammar kind.}
\label{tab:node_dist}
\begin{tabular}{lr@{\hskip 2em}lr}
\toprule
Type & Share & Type & Share \\
\midrule
$\bot$ (\texttt{UNKNOWN}) & 75.23\% & \texttt{COMPARISON} & 2.72\% \\
\texttt{PTR\_DEREF} & 6.56\% & \texttt{ARITH\_OP} & 1.57\% \\
\texttt{CALL\_SITE} & 5.45\% & \texttt{ARRAY\_INDEX} & 0.53\% \\
\texttt{CONTROL\_BRANCH} & 4.18\% & \texttt{MEMORY\_ACCESS} & 0.30\% \\
\texttt{ASSIGN} & 3.47\% &  &  \\
\bottomrule
\end{tabular}
\end{table}


\paragraph{The data are strongly imbalanced, and we do not rebalance them.}
Only \numbigvulPrev{} of BigVul and \numdiversevulPrev{} of DiverseVul
functions are labelled vulnerable. We handle this in three ways. (i)~Training
uses a class-weighted loss, with positive weight equal to the training
benign-to-vulnerable ratio capped at 20. (ii)~Test sets are never resampled, so
every metric reflects the true prevalence. (iii)~We report AUC and AUPRC, which
do not depend on a decision threshold, together with the trivial
``everything is vulnerable'' classifier as a floor; its $F_1$ at the observed
test prevalence is only $2p/(1{+}p)$.

\paragraph{Why overall prevalence and test prevalence differ.} In BigVul the
overall prevalence is \numbigvulPrev{} but the pooled prevalence of the
held-out folds is \numbigvulTestPrev{}. The reason is the train-only rule of
Section~\ref{sec:protocol}: the largest repositories (Chrome and the Linux
kernel; \numbigvulTrainOnlyN{} functions, prevalence \numbigvulTrainOnlyPrev{})
are never tested, and the remaining projects have a higher vulnerable share. An
earlier version of this paper quoted these two figures (0.054 and 8.8\%) without
connecting them.

\paragraph{Project concentration.} Even after the largest repositories are held
train-only, one project can dominate a test fold: in BigVul the repository
\texttt{\numbigvulTopTestProject} makes up as much as
\numbigvulTopTestShareMax\% of a fold's test functions (the largest share in any
DiverseVul fold is \numdiversevulTopTestShareMax\%). We therefore report every
main BigVul comparison a second time with that project removed from the test
predictions (Tables~\ref{tab:rq2_representation} and~\ref{tab:perproject} and
the text of Section~\ref{sec:rq_representation}).

\paragraph{The 200-node cap.} A graph keeps the first 200 named nodes in
pre-order. Table~\ref{tab:corpora} shows that this truncates
\numbigvulTruncPct\% of BigVul functions and \numdiversevulTruncPct\% of
DiverseVul functions, and that it is far from class-neutral: in BigVul
\numbigvulTruncVulnPct\% of vulnerable functions are truncated against
\numbigvulTruncBenignPct\% of benign ones (Fisher exact test
$p$~\numbigvulTruncP). Vulnerable functions are simply longer. This matters in
two ways, both examined below: a model may learn function length rather than
anything about vulnerabilities (Section~\ref{sec:rq_reference}), and what lies
beyond the cap is invisible to every compared system alike.

\subsection{Deduplication}
\label{sec:dedup}
Near-duplicate functions are removed \emph{before} any split. Vendored code
appears under several project names, so without this step a held-out function
can have a twin in training. Signatures are computed over the abstracted node
sequences ($\kappa$ and $\phi$), which catches renamed copies that text hashing
misses; this removes 2--6\% of the sampled functions (the remaining counts are
in Table~\ref{tab:corpora}).

\subsection{Cross-project protocol}
\label{sec:protocol}
Evaluation uses project-level GroupKFold with ten folds in every experiment:
projects are dealt into five folds twice, with two different shuffles, and no
project is ever split between training and test. Projects too large to be placed
without dominating a fold (those larger than 30\% of the mean fold size: Chrome and the Linux kernel in
BigVul, the Linux kernel only in DiverseVul) are held \emph{train-only}. Every
condition in every table is run on the same ten folds; the analysis scripts
refuse to merge result files whose fold assignments differ. The two repeats are
\emph{not} independent: the splitter deals the projects out largest first, so
the largest held-out projects fall into the same fold index in both repeats
(for example the repository that dominates BigVul test folds appears in folds 0
and 5), while the small projects are reshuffled (the test project sets of fold
$k$ and $k+5$ overlap by only 13--23\% in Jaccard index). The ten folds are
therefore better described as five folds with partial replication.
Section~\ref{sec:stats} says how the tests account for this.

\paragraph{Per-fold training.} The lexical dictionary is built from the training
projects only (Section~\ref{sec:features}); the classifier is initialised and
shuffled with seed $1000+k$ on fold $k$, the same seed for every compared
system. Hyper-parameters are in Table~\ref{tab:hparams}. They were fixed before
the experiments and are not tuned on any fold; there is no validation split and
no early stopping, so every system trains for exactly \numcfgEpochs{} epochs.

% Generated by experiments/emit_tables.py -- DO NOT EDIT BY HAND.
% Every value is read from experiments/results/*.json.

\begin{table}[t]
\centering
\caption{Hyper-parameters. All values are fixed in advance and shared by every
compared system; none is tuned on any fold.}
\label{tab:hparams}
\begin{tabular}{ll}
\toprule
Item & Value \\
\midrule
Embedding width / hidden width & 64 / 128 \\
Message-passing layers & 2 (GGNN: 2 gated steps) \\
Transformer & 2 layers, 4 heads, feed-forward $2\times$ width, dropout 0.1 \\
Readout and head & mean $\Vert$ max pooling; MLP $2h\to h/2\to 1$, dropout 0.3 \\
Optimiser & Adam, learning rate 0.001 \\
Batch size & 64 \\
Epochs & 10 (fixed; no early stopping, no validation split) \\
Loss & binary cross-entropy, positive weight $=\min(n_{\text{neg}}/n_{\text{pos}}, 20)$ \\
Seeds & initialisation and shuffling seed $1000+k$ for fold $k$, shared by all systems \\
Lexical dictionary & 10{,}000 most frequent training-fold tokens with count $\ge 2$ \\
Node cap & 200 named AST nodes (pre-order) \\
Sample & 10,000 functions per corpus, uniform at random, seed 0 \\
Folds & 5$\times$2 = 10 project-level folds \\
Pretrained reference & UniXcoder-base, 256 tokens, frozen; logistic regression $C=0.1$, balanced \\
\bottomrule
\end{tabular}
\end{table}


\subsection{Statistical treatment}
\label{sec:stats}
Every paired comparison is made between two systems on the folds both
completed. The reported difference is the \emph{mean over folds} of the
per-fold difference in the metric, which is the quantity the per-fold columns of
the tables average. We attach three kinds of uncertainty statement.
\begin{enumerate}
    \item A \textbf{cluster bootstrap over held-out projects}. In each of the
          resamples, projects are drawn with replacement (a project is one
          cluster however many folds it is tested in); every function of a drawn
          project is repeated as often as the project was drawn; the metric
          difference is recomputed inside each fold and averaged. The 95\%
          percentile interval and a two-sided $p$ follow. Functions inside a
          repository share style, APIs and often code, so resampling functions
          instead would understate the uncertainty. We use 1{,}000 resamples
          under each of five seeds for AUC (the verdict is called stable if the
          interval excludes zero under at least 90\% of the seeds), and 500
          resamples under one seed for AUPRC and the size-binned AUC.
    \item A \textbf{Wilcoxon signed-rank test}~\cite{wilcoxon1945individual}, a
          paired test that uses only the ranks of the ten per-fold differences
          (exact, two-sided; its smallest attainable $p$ is $0.002$). Because
          fold $k$ and fold $k+5$ are partial replicates of each other
          (Section~\ref{sec:protocol}), we also run it on the five
          repeat-averaged differences, where the smallest attainable $p$ is
          $0.0625$, so that test can never reach the 5\% level on its own.
    \item \textbf{Holm-adjusted $p$-values}~\cite{holm1979simple}, which
          correct a family of tests for the chance that one of them looks
          significant by luck. The family is all contrasts for one corpus and
          metric; adjusted values are stored in the result files and quoted where
          a conclusion depends on one contrast.
\end{enumerate}
We therefore treat the cluster bootstrap as the primary inference, the Wilcoxon
test as supporting evidence, and call a difference supported only if the
bootstrap interval excludes zero. Whether truncation depends on the label is a
question about counts, for which we use Fisher's exact test (a small-sample
alternative to the chi-square test of independence).

\subsection{Compute and reproducibility}
All experiments run on CPU. One run of the BigVul GGNN grid takes roughly
seven minutes per cell per core. Every table and every number in the text is
generated from the stored per-fold predictions by scripts in the repository
(\texttt{experiments/run\_all.sh}); no result is typed by hand.
